%%%%%%%%%%%%%%%%%%%%%%%%%%%%%%%%%
%
%       EXERCÍCIO
%
%%%%%%%%%%%%%%%%%%%%%%%%%%%%%%%%%

\ifdefstring{\atividade}{prova}{%
    \ifdefstring{\modo}{objetivo}{%
        \renewcommand{\valorquestao}{\ValorQObj\ ponto}
    }{%
        \renewcommand{\valorquestao}{\ValorQDisc\ pontos}
    }%
}%

\begin{exercicioBanco}[\valorquestao]
Considere as afirmações a seguir.
\begin{enumerate}[I.]
    \item A função afim \(f(x) = -x + 2\) é injetora.
    %
    \item Seja \(f: \bbR \to \bbR\) um função afim com coeficiente angular negativo e tal que \(f(2) = 0\). Então \(f(x) < 0\), para todo \(x > 2\).
    %
    \item O perímetro de um retângulo com medidas \(10\) cm e \(x\) cm é maior do que \(30\), se \(x \ge 5\).
\end{enumerate}
%
Assinale a alternativa correta.

% Define as alternativas
\newcommand{\alternativas}{%
    \begin{center}
        \begin{tabularx}{\textwidth}{XXX}
            (a) Apenas I é verdadeira. &
            (b) I e II são verdadeiras. &
            (c) Apenas III é verdadeira. \\[5pt]
            (d) I e III são verdadeiras. &
            (e) Todas são verdadeiras.
        \end{tabularx}
    \end{center}
}

% Resposta correta
\newcommand{\resposta}{B}

% Lógica condicional para exibição
\ifdefstring{\atividade}{lista}{%
    \alternativas
    \vspace{0.5em}
    
    \noindent\textbf{Resposta:} letra \textbf{\resposta}.
}{%
    \ifdefstring{\modo}{objetivo}{%
        \alternativas
    }{%
        % modo = discursiva → não mostra alternativas
    }
}
\end{exercicioBanco}

